\section{Discussion}\label{sec:discussion}

The main advance over the GA-based precursor is methodological control. The earlier work offered a useful proof of concept for combining technical indicators, feature selection, and GA tuning, but it could not distinguish optimizer behaviour from the consequences of randomized validation and unequal search opportunity. The present design makes these assumptions visible: chronological data ordering is preserved, evaluation budgets are fixed, and the validation criterion is stated separately from the final test metrics.

The direct comparison does not support a simple rule that accuracy fitness should dominate when accuracy is reported, or that MCC/$F_1$ fitness should dominate every robust metric. The paired cells show positive and negative changes under each criterion. This is plausible in intraday markets, where validation performance is affected by small decision-boundary changes and the final chronological block may represent a different regime. For this reason, the appropriate interpretation is objective-dependent behaviour, not a universal ranking of the two fitness functions.

The within-protocol exploratory analysis provides a more bounded result. MLP ranks first under the accuracy-fitness protocol across matched optimizer--seed blocks, and its adjusted pairwise contrasts against the other backbones are significant. The economic analysis, however, demonstrates that the predictive ordering does not settle the decision alone: RF attains the smallest mean drawdown under accuracy fitness, whereas MLP obtains the largest mean profit. A deployment decision must therefore specify whether return, downside risk, turnover, or stability is the governing criterion.

Several limitations should guide interpretation. Only three seeds are common to all model--optimizer cells in both protocols. The direct dual-fitness results are therefore descriptive, and the model-level Friedman/Wilcoxon analyses are exploratory. The data cover one liquid Brazilian futures market and one chronological test segment; results may differ across instruments and regimes. Finally, the fixed 1,000-evaluation allowance is appropriate for a controlled benchmark but does not reveal long-budget asymptotic behaviour.

The next study phase is clear: run a preregistered common 20-seed design for every cell, retain the same temporal split and candidate budget, record all validation metrics at each candidate evaluation, and apply paired tests separately to MCC, accuracy, and economic outcomes. That extension would permit confirmatory inference about fitness choice while preserving the methodological strengths established here.
